\chapter{Systemic Lupus Erythematosus}
\index{Systemic Lupus Erythematosus}
\label{sec:Systemic Lupus Erythematosus}

\section{Overview}
Systemic lupus erythematosus is a chronic, multisystem autoimmune disease characterized by loss of tolerance to self-antigens and production of a wide array of autoantibodies, leading to immune complex deposition and inflammation in multiple organs, predominantly affecting women of childbearing age (9:1 female predominance) and being more common in African Americans, Hispanics, and Asians. This autoimmune condition occurs when the immune system mistakenly attacks the body's own tissues. It follows a chronic course with periods of flare and remission. Autoimmune diseases result from an abnormal immune response where the body mistakenly attacks its own healthy tissues. These conditions are typically chronic, with alternating periods of flare and remission. They can affect a single organ or be systemic, involving multiple organ systems. Women are affected more frequently than men for most autoimmune conditions. The exact prevalence varies by condition, but collectively autoimmune diseases affect a significant portion of the population.
\section{Causes}
This condition happens because of autoantibody production and immune complex deposition. In autoimmune conditions, a combination of genetic susceptibility and environmental triggers initiates an inappropriate immune response against self-antigens. This leads to chronic inflammation and tissue damage in affected organs. A combination of genetic susceptibility and environmental triggers initiates the autoimmune process. Specific HLA genotypes are strongly associated with many autoimmune conditions. Environmental triggers may include infections, smoking, medications, and hormonal changes. Loss of self-tolerance leads to activation of autoreactive T cells and B cells producing autoantibodies. Molecular mimicry between pathogen antigens and self-antigens may trigger cross-reactive immune responses.
\section{Risk Factors}


\begin{itemize}[noitemsep]
  \item Family history of autoimmune disease
  \item Female sex (higher prevalence)
  \item Specific HLA genotypes
  \item Epstein-Barr virus infection
  \item Cigarette smoking
  \item Certain medications (drug-induced autoimmunity)
  \item Hormonal factors (pregnancy, postpartum)
  \item Vitamin D deficiency
\end{itemize}
\section{Signs and Symptoms}

\subsection{Common symptoms}
\begin{itemize}[noitemsep]
  \item Clinical features are protean and include malar rash, photosensitivity, oral ulcers, arthritis, serositis (pleuritis, pericarditis), nephritis (lupus nephritis---the most serious manifestation), cytopenias, seizures, and psychosis. Fatigue and general malaise Joint pain, swelling, and stiffness Skin rashes and lesions Low-grade fever Muscle weakness and pain Organ-specific symptoms based on affected system Weight changes (loss or gain) Dry eyes and dry mouth (Sjogren syndrome) Raynaud phenomenon (cold-induced color changes in fingers) Fluctuating symptoms with periods of flare and remission
  \item Fatigue and general malaise
  \item Joint pain, swelling, and stiffness
  \item Skin rashes and lesions
  \item Low-grade fever
  \item Muscle weakness and pain
  \item Organ-specific symptoms based on affected system
  \item Weight changes (loss or gain)
  \item Dry eyes and dry mouth (Sjogren syndrome)
  \item Raynaud phenomenon (cold-induced color changes in fingers)
  \item Fluctuating symptoms with periods of flare and remission
\end{itemize}

\subsection{Emergency warning signs}
\begin{itemize}[noitemsep]
  \item Severe or worsening symptoms of systemic lupus erythematosus require immediate medical evaluation
\end{itemize}
\section{Progression and Stages}
The 2019 EULAR/ACR classification criteria require an ANA titer $\geq$1:80 plus weighted criteria across clinical and immunological domains. Autoimmune diseases typically follow a relapsing-remitting course with unpredictable flares. Early disease may present with nonspecific symptoms before organ-specific manifestations develop. Chronic inflammation over time can lead to irreversible tissue damage and organ dysfunction. With appropriate treatment, many patients achieve remission or low disease activity states.
\section{Complications}
\begin{itemize}[noitemsep]
  \item Irreversible organ damage from chronic inflammation
  \item Joint deformities and erosions in inflammatory arthritis
  \item Renal failure in lupus nephritis
  \item Pulmonary fibrosis or hypertension
  \item Cardiovascular complications from systemic inflammation
  \item Osteoporosis from chronic corticosteroid use
  \item Increased risk of infections from immunosuppressive therapy
  \item Lymphoma risk in some autoimmune conditions
\end{itemize}
\section{Diagnosis}
Diagnosis typically involves antinuclear antibody (ANA, sensitive but not specific), anti-dsDNA (specific, correlates with nephritis), anti-Smith (highly specific), complement levels (C3, C4 decreased during flares), and kidney biopsy for lupus nephritis classification (ISN/RPS). Diagnosis is based on a combination of clinical features, laboratory findings (autoantibodies, inflammatory markers), and imaging studies. Classification criteria help standardize the diagnostic process. Clinical history and physical examination Autoantibody testing (ANA, RF, anti-CCP, anti-dsDNA, etc.) Inflammatory markers (ESR, CRP) Complete blood count and comprehensive metabolic panel Imaging studies (X-ray, MRI, ultrasound) of affected areas Biopsy of affected tissue for histological examination Classification criteria for specific autoimmune diseases Rheumatology consultation for suspected autoimmune disease Monitoring for organ involvement (renal, pulmonary, cardiac) Genetic testing for HLA associations when indicated

\begin{itemize}[noitemsep]
  \item Clinical history and physical examination
  \item Autoantibody testing (ANA, RF, anti-CCP, anti-dsDNA, etc.)
  \item Inflammatory markers (ESR, CRP)
  \item Complete blood count and comprehensive metabolic panel
  \item Imaging studies (X-ray, MRI, ultrasound) of affected areas
  \item Biopsy of affected tissue for histological examination
  \item Classification criteria for specific autoimmune diseases
  \item Rheumatology consultation for suspected autoimmune disease
\end{itemize}
\section{Prevention}
\begin{itemize}[noitemsep]
  \item No known prevention for most autoimmune diseases
  \item Avoid known triggers when possible
  \item Maintain a healthy lifestyle to support immune function
  \item Early recognition of symptoms for prompt diagnosis
  \item Regular medical check-ups for monitoring
  \item Adequate vitamin D levels through sun exposure or supplementation
\end{itemize}
\section{Treatment Options}
Treatment options include hydroxychloroquine (cornerstone for all patients, reduces flares and mortality), corticosteroids, and immunosuppressive agents (mycophenolate, azathioprine, cyclophosphamide for severe organ involvement), with newer targeted therapies including belimumab (anti-BLys/BAFF monoclonal antibody), voclosporin, and anifrolumab (anti-type I interferon receptor). Treatment aims to control inflammation, manage symptoms, and prevent disease flares. A step-up approach is often used, starting with symptomatic treatments and progressing to immunosuppressive therapies as needed. Nonsteroidal anti-inflammatory drugs (NSAIDs) for symptom relief Corticosteroids for rapid control of inflammation Disease-modifying antirheumatic drugs (DMARDs) Biologic therapies targeting specific immune pathways Immunosuppressive agents (azathioprine, mycophenolate, cyclophosphamide) Antimalarial drugs (hydroxychloroquine) for mild disease Physical and occupational therapy to maintain function Pain management for chronic pain Lifestyle modifications including diet and stress reduction Regular monitoring for disease activity and treatment side effects

\begin{itemize}[noitemsep]
  \item Nonsteroidal anti-inflammatory drugs (NSAIDs) for symptom relief
  \item Corticosteroids for rapid control of inflammation
  \item Disease-modifying antirheumatic drugs (DMARDs)
  \item Biologic therapies targeting specific immune pathways
  \item Immunosuppressive agents (azathioprine, mycophenolate, cyclophosphamide)
  \item Antimalarial drugs (hydroxychloroquine) for mild disease
  \item Physical and occupational therapy to maintain function
  \item Pain management for chronic pain
  \item Lifestyle modifications including diet and stress reduction
\end{itemize}
\section{Living with It}
\begin{itemize}[noitemsep]
  \item Take medications consistently as prescribed
  \item Identify and avoid personal flare triggers
  \item Maintain a balanced diet and regular gentle exercise
  \item Practice stress management techniques
  \item Join support groups for emotional support
  \item Work with a rheumatologist for ongoing disease management
  \item Monitor for medication side effects
  \item Plan activities around energy levels during flares
\end{itemize}
\section{Outlook}
The outlook varies from person to person, with improved outcomes due to early diagnosis and better management. Autoimmune conditions are typically chronic and require long-term management. With appropriate treatment, many patients achieve good symptom control and maintain quality of life. Autoimmune diseases are generally chronic conditions requiring long-term management. Disease activity can fluctuate, requiring adjustments in treatment over time. Early diagnosis and treatment improve outcomes and prevent irreversible organ damage. Research continues to develop more targeted therapies with fewer side effects.
